% Generated by `taskgen course examples` from dlp.brute_force (seed 4).
% Do not edit: change examples.yaml and regenerate.
\begin{example}[перебор, $h \notin \langle g \rangle$]\label{ex:dlp-brute-force}
Решить задачу дискретного логарифмирования перебором, вычисляя последовательно $g, g^2, g^3, \dots$: \[2^x \equiv 14\ (\textrm{mod}\ 23).\] Найти все решения или показать, что их нет. Сколько умножений потребовалось?

\par\noindent\textit{Решение.} Каждая следующая степень --- одно умножение на $g = 2$ по модулю $23$; идём до возврата в $1$ (тогда $k = \ord(g)$ и степени начинают повторяться):\[\begin{array}{c|rrrrrrrrrrr}k & 1 & 2 & 3 & 4 & 5 & 6 & 7 & 8 & 9 & 10 & 11 \\ \hline g^k & 2 & 4 & 8 & 16 & 9 & 18 & 13 & 3 & 6 & 12 & 1\end{array}\]Итак, $n = \ord(2) = 11$, $\langle g \rangle$ --- подгруппа индекса $2$ в $(\mathbb{Z}/23\mathbb{Z})^*$. Число $14$ в списке $\langle g \rangle$ не встретилось, а дальше степени повторяются ($g^{k + 11} = g^k$); значит, $h \notin \langle g \rangle$ (понадобилось $10$ умножений). Проверка: элементы $\langle g \rangle$ --- в точности решения $y^{11} = 1$, а $14^{11} \equiv 22 \ne 1$. \[\text{Решений нет.}\]
\end{example}
